% !TeX root = ../papers/nc-p3.tex
\section{点と二平面の巡回からの全次数判定}
\label{sec:point-line-flag-cycles}
四線型テンソル $w$ の隣接二対二切断がともに階数6、一文字切断が全階数とし、
各辺の六次元像 $R_i\subset V_i\otimes V_{i+1}$ で定まる二次代数を $A$ とする。
両向きの非零評価が全て階数3以上で、どの $R_i$ も階数1テンソルを含まないとする。
以下の点は関係で結ばれた両無限点列に属するものを使う。
位置 $i$ の点 $p$ の消滅超平面を $H(p)$、前点を $p^-$ と書く。
$U\subset H(p)$ が二平面なら、$U^\perp$ で $p$ から独立な $q$ を一つ取り、
\[
 E(U,p)=\operatorname{supp}_{\rm left}w(p^-,q,-,-)\ \subset V_{i+1}
\]
とする。$w(p^-,p,-,-)=0$ のため選び方に依存しない。

\begin{theorem}\label{thm:recurrent-point-line-flags}
各位置で有限個の二平面を用意し、以下の二条件を満たすとする。有限周期に展開してよい。
\begin{enumerate}
\item 各二平面 $U$ に $U\subset H(p)$ となる指定点があり、$E(U,p)$ は次の位置で
用意した二平面の一つである。
\item この集合から各位置で $U_i$ を選んだ周期列があり、
$R_i\cap(U_i\otimes U_{i+1})\ne0$ が全ての位置で成り立つ。
\end{enumerate}
このとき $A$ は全次数で標準Koszul左右AS正則である。
有限周期の指定点列と二平面の遷移先を固定したこの判定は、上の評価下界と階数1関係の不在を
保つ特殊化の下で閉じている。
\end{theorem}
\begin{proof}
点 $p$ の右点加群を $M(p)$ とする。$H(p)$ の右イデアルはその全関係イデアルで、
$H(p)\otimes e_{i+1}A\to e_i A$ の核は
$K(p)=R_i\cap(H(p)\otimes V_{i+1})$ の三元で生成される。
命題\ref{prop:commuting-three-point-graph}の点上のTor計算を用いるが、
ここで必要な部分は標準有限双対をまだ仮定せずに通る。
実際、終点単純の左極小解消の最初の三項は矢と二次関係で決まり、次の項に
$w$ の一文字縮約の四次元空間を三次の極小重なりとして含められる。
点加群にテンソルしたこれらの射の階数は内部次数3以上で $1,3,3$ である。
二番目は関係の点評価、三番目も反対の隣接切断の完全対合を通した点評価の転置である。
従って第一Torは次数1で三個、第二Torは次数2の三個の商となる。
実際に $K(p)$ の三元が独立な二次関係を与える。

$h=w(p^-,-,-,-)$ の最終文字の縮約は $K(p)$ を張る。
その像は $K(p)$ に含まれ、最終文字側での階数3は $p^-$ への前辺の評価下界と
相補切断の完全対合から従う。$q$ で $K(p)$ を評価した像は $E(U,p)$ である。
$L_i(U)=e_i A/UA$ とし、$H(p)/U$ の生成元で $q$ が1を取る持上げを $v$ とする。
上記の三関係を $q$ で読むと $(UA:v)=E(U,p) A$、従って
\[
 0\longrightarrow L_{i+1}(E(U,p))(-1)\xrightarrow{\ v\ }L_i(U)
 \longrightarrow M(p)\longrightarrow0 .
\]
任意の長さについて条件1を反復すると $\dim L_i(U)_n=n+1$ である。
さらに $UA$ の関係はただ一つの二次関係 $r\in R_i\cap(U\otimes V_{i+1})$ で生成される。
実際 $K(p)$ の $q$ 評価の核は $kr$ であり、その基底を $r,s_1,s_2$ に延ばす。
$h=r\otimes c_0+s_1\otimes c_1+s_2\otimes c_2$ とすると
$r'=q(s_1)\otimes c_1+q(s_2)\otimes c_2$ が次の平面の非零関係である。
非零性は $p^-$ の唯一の後続が $p$ であることによる。
所与の $U$ の関係を $ra_0+s_1a_1+s_2a_2$ と表し、一つ短い長さで帰納仮定を
$E(U,p)$ に使うと $(a_1,a_2)=(c_1b,c_2b)$ となる。
$h\in H(p)\otimes R_{i+1}$ だから元の関係は $r(a_0-c_0b)$ である。

条件2の各 $r_i$ は純関係がないので階数2である。$U_i$ と $kr_i$ による抽象的二文字代数 $P$ は
位置ごとの基底で通常の交換子表示になり、長さ $n$ で次元 $n+1$ を持つ。
$L_i(U_i)$ の基底を持ち上げ左から $P$ の語を掛けると $A$ を張るので、各成分のHilbert上界は
$(1-z)^{-4}$ である。固定終点の経路次元を $a_n$ とすると上の一関係表示から
$a_n-2a_{n-1}+a_{n-2}=n+1+e_n$, $e_n\ge0$ を得る。
$(1-z)^{-2}$ の係数により上界で等号である。従って持上げは左 $P$ 自由基底を与え、
どの $U_i$ の任意の非零文字による左乗法も単射である。

任意の両無限点列について、点加群の右表示を三次の列 $h$ で延ばす。
得られる線形列の階数は $1,3,3,1$ で、既に項1まで完全である。
上の $c_0,c_1,c_2$（一般の点では $K(p)$ の任意の基底についての三係数）は独立である。
この三平面は同じ位置の $U_{i+2}$ と交わるので、三次の列の少なくとも一つの線型結合が
単射な左乗法となり、最後の列は単射である。点加群のHilbert関数と得られた標準Hilbert関数から
中央も完全となる。従って点加群とその次の点加群はともに線形解消を持つ。
後者を一つ移動して前者に入れた商は始点単純なので、錐によるその解消も線形である。
全始点でのKoszul性から双対の次元は $1,4,6,4,1$、$w$ とその縮約は全重なりとなる。
一文字の全階数と隣接切断の階数6で相補次数の積対合は完全なので、左右AS条件が従う。

最後に曲線上の付値環で、一般側の点、全平面、条件2の関係直線を各射影空間・Grassmann多様体に延ばす。
点列と関係の包含は閉条件である。$U\subset H(p)$ で独立な $q$ を特殊側で選び直し、
局所自由な $U^\perp$ に持ち上げる。一般側で $w(p^-,q,-,-)$ の左支持が次の平面に入るので
特殊側でも入る。$q\ne p$ と一意な隣接点から非零で、純関係がないため階数2のままである。
全ての遷移が残り、条件2の非零関係直線の特殊化も階数2を保つ。付値判定または射影像の閉性で従う。
\end{proof}

\begin{corollary}\label{cor:rank-ten-sheets-complete-regularity}
節\ref{sec:equal-determinant-second-sheet}の $Dpw\ne0$ の第二補完では
$(a^2-b^2)(d^2-c^2)(a^2-c^2)(d^2-b^2)\ne0$ が標準Koszul左右AS正則性の必要十分条件である。
節\ref{sec:unequal-determinant-conic}の $D\beta(1-\beta)\ne0$ の円錐曲線族では
$(1+\beta)I J K N_0\ne0$ が必要十分である。ともに三重衝突を含む。
従って $z=z\lambda^2=-1$ の可逆な薄い繊維では、全表示の各閉点の正則性が決定される。
\end{corollary}
\begin{proof}
必要性は両節の評価障害である。両厚色で関係空間は低位と高位への射影がともに同型であり、
各薄い色は全成分が非零である。階数1の関係があれば組をまたぐ二部分行列の行列式はともに零となる。
二薄色の係数を $x,y$ とすると、符号・単元を除いて $x^2-y^2$ と $x^2+y^2$ なので両係数が零である。
残った対角二ブロックの片方は零となり、二厚色の同型射影によりもう片方も零である。
逆辺にも同じ計算が使える。

\path{computations/verify_point_line_flag_cycles.py}は両枝の関数体上で以下の有限証明書を検算する。
等行列式では $b=B^2,c=C^2$ として位置1の表から五点を取る。
点名 $1,4,9,10,15$ は、$a=2,B^2=3,C^2=5,d=7$ の素数109での座標が順に
$(0,0,1,1),(0,1,47,26),(1,0,59,10),(1,0,50,99),(0,1,62,83)$ となる符号を指定するだけである。
他の位置の点は四文字縮約の一意な核へ順に移す。
位置0,2で、前々位置の点4と前位置の点1による縮約の左支持を $U_i$、右支持を $U_{i+1}$ に取る。
位置1,3では同じ五点の収縮のうち両支持が所定の二平面となるものを取れる。
これが条件2である。条件1の指定点を、各 $U_i$ から順に
\[
 (1,9,10,15,4),\quad(15,4,10,9),\quad(1,9,10,15,4),\quad(15,4,9,10)
\]
と選ぶ。五歩列では最後が一歩後の平面に戻り、四歩列では最初に戻る。
従って平面の和集合に必要な遷移が存在する。計算は各 $4\times4$ 縮約の二行の外積と
三行の外積で階数と平面の一致を整数多項式として確認する。

円錐曲線族では共通倍率でまず $b=1$ とし、$l^2=-I/J$ と取れば
$e=(kd-1)/(d-k)$, $a=((e-1)-l^2(e+1))/((e-1)+l^2(e+1))$ として
$k,d,l$ が独立なパラメータになる。位置0で二つの点対を使う。
第一対は $v=0$ 面の $(1,0,j,lj),(1,0,-j,-lj)$、第二対は $u=0$ 面の
$(0,1,v,Lv),(0,1,-v,-Lv)$ とし $L=-(k-1)/((k+1)l)$ である。
$j^2=J/(2(a+e)),v^2=N_0/(2(d+h))$ である。位置別に各対の割線の消滅二平面を作る。
偶では第一対、奇では第二対を取ると、各辺で六関係にこの二平面の積四式を加えた階数は9である。
従って条件2が成立する。どの位置のどちらの割線からも、その二点を順に使う二歩で
二位置後の他方の割線の消滅二平面に移る。これを繰り返して条件1を得る。
この二歩の検査は、使う対ごとに共通倍率で $j=1$ または $v=1$ として行える。
他方の対の二平面は二点の低位部と高位部の各直線で決まり、平方根を要しない。
従ってここでも三独立変数の多項式恒等式で十分である。
二つの既約表示に定理の特殊化を適用すれば評価障害の補集合の全閉点が正則となる。
最後に等行列式の $p w=0$ と正負の交換表示は既出の全点判定を持ち、その他の表示は
節\ref{sec:full-thin-parallel-cover}の移動・捻りで以上に戻る。
\end{proof}
